%% ma: L10-B15-0086
%% lop: 10
%% chuong: 6
%% bai: 15
%% dang: 10.15.6
%% loai: TLN
%% muc: VD
%% dap_an: "353"
%% nguon: "Ảnh giáo viên gửi 25/11/2026: Bài tập trắc nghiệm – Toán 10 Bài 15 Hàm số (3 phần), Phần III Câu 3"
%% kien_thuc: "Lớp 10 Bài 15 (hàm số, tập giá trị trên một đoạn; hàm bậc nhất đồng biến)"
%% trang_thai: chinh-thuc
\begin{ex}
Một người cần đặt một tiệc cưới ước tính khoảng $30$ đến $35$ bàn. Nhà hàng thứ nhất đề nghị anh nay đóng tiền cố định $20$ triệu đồng, sau khi tiệc cưới diễn ra sẽ đóng khoản còn lại với số tiền $2$ triệu đồng/1 bàn. Nhà hàng thứ hai đề nghị anh đóng tiền cố định $10$ triệu đồng, sau khi tiệc cưới diễn ra sẽ đóng khoản còn lại với số tiền $2{,}5$ triệu/1 bàn. Nếu anh này nên lựa chọn nhà hàng thứ nhất thì phải trả số tiền trong $[a;b]$, $a,b>0$, $a,b$ có đơn vị triệu đồng. Nếu anh này chọn nhà hàng thứ hai thì phải trả số tiền trong $[c;d]$, $c,d>0$, $c,d$ có đơn vị là triệu đồng. Tính $a+b+c+d$ (kết quả làm tròn đến hàng đơn vị).
\shortans{353}
\loigiai{Với $x\in[30;35]$ bàn: nhà hàng thứ nhất $f(x)=20+2x\in[80;90]$; nhà hàng thứ hai $g(x)=10+2{,}5x\in[85;97{,}5]$. $a+b+c+d=352{,}5\approx353$.}
\end{ex}
